\begin{longtable}{TD}
\caption{Task behavior terms.}
\label{tab:terms-task-behavior}\\

\toprule
Term & Definition \\
\midrule
\endfirsthead

\toprule
Term & Definition \\
\midrule
\endhead

\bottomrule
\endfoot

\termrow{behavior}{
The observed or estimated execution pattern of a task kind, task instance, or
task attempt under a particular execution context. Behavior includes resource
demand, worker occupancy, retry behavior, dependency wait, outcome,
uncertainty, and risk-relevant effects.
}

\termrow{task behavior}{
The way a task executes: which resources it demands, which bounded capacities
it holds, how it completes, how variable it is, and what operational
consequences may follow from misestimating or mishandling it.
}

\termrow{behavior profile}{
A summarized description of how a task kind behaves under a class of execution
contexts. A behavior profile may include demand distributions, occupancy
distributions, dominant behavior classes, uncertainty, risk class, and relevant
features.
}

\termrow{behavior class}{
An analytical category derived from observed or estimated behavior. Examples
include compute-dominant, dependency-bound, retry-amplifying, worker-slot-heavy,
payload-sensitive, state-dependent, or high-tail-risk behavior.
}

\termrow{behavior mode}{
A context-dependent mode of behavior for a task kind. The term emphasizes that
the same task kind may behave differently under different execution contexts.
}

\termrow{compute-dominant behavior}{
Behavior in which CPU demand is the dominant active resource dimension. A task
kind may exhibit compute-dominant behavior without being intrinsically a CPU
task in all contexts.
}

\termrow{memory-pressure-sensitive behavior}{
Behavior that is sensitive to memory pressure, allocation rate, working-set
size, or garbage-collection effects. This behavior is not necessarily identical
to being memory-bound.
}

\termrow{disk-IO-dominant behavior}{
Behavior in which disk reads, disk writes, storage throughput, or I/O wait
dominates execution. Disk I/O demand is distinct from disk pressure in the
execution environment.
}

\termrow{network-sensitive behavior}{
Behavior that changes substantially with network condition, such as packet
loss, jitter, retransmissions, bandwidth changes, or timeouts.
}

\termrow{dependency-bound behavior}{
Behavior constrained primarily by external dependencies, such as databases,
APIs, remote services, rate limits, quotas, or connection pools.
}

\termrow{retry-amplifying behavior}{
Behavior in which retries substantially increase work, worker occupancy, or
operational risk. Retry count may be a symptom or outcome rather than an
independent cause.
}

\termrow{worker-slot-heavy behavior}{
Behavior with high worker-slot time relative to active resource usage. This
class is important for queue workers because a task may consume little CPU but
hold scarce worker capacity for a long time.
}

\termrow{payload-sensitive behavior}{
Behavior that depends strongly on payload size, input shape, record count,
object count, or batch size. Payload-sensitive behavior can still be uncertain
if context or hidden state varies.
}

\termrow{state-dependent behavior}{
Behavior that depends on domain-specific or internal system state, such as
replica divergence, cache state, compaction debt, segment fragmentation, or
dirty data volume.
}

\termrow{high-tail-risk behavior}{
Behavior in which rare slow executions or high-percentile outcomes are
operationally important. High-tail-risk behavior is distinct from high average
cost.
}

\termrow{stable behavior}{
Behavior with low residual error, low drift, and low variability under
comparable execution contexts. Stable behavior may still be expensive.
}

\termrow{unstable behavior}{
Behavior with high variability, drift, residual error, multimodality, or
context sensitivity. Unstable behavior is not necessarily high-cost in every
attempt.
}

\end{longtable}
